% TOP_PROBLEM: {"id": 12, "title": "Goldbach's Conjecture", "category_id": 1, "category": {"id": 1, "name": "number_theory", "display_name": "Number Theory", "description": "Properties of integers, prime numbers, Diophantine equations.", "slug": "number-theory", "order_index": 1, "created_at": "2026-07-31T15:26:25.670Z"}, "status": "open"}
% FIRST_POSTED: 2026-09-27
% !TeX program = lualatex
% Standalone notebook; compile twice: lualatex 12.tex
% Original section preserved verbatim except for marked insertion blocks.
\documentclass[11pt,a4paper]{article}
\usepackage[margin=24mm,headheight=14pt]{geometry}
\usepackage{fontspec}
\setmainfont{Latin Modern Roman}
\setsansfont{Latin Modern Sans}
\setmonofont{Latin Modern Mono}
\usepackage{amsmath,amssymb,mathtools}
\usepackage{unicode-math}
\setmathfont{Latin Modern Math}
\usepackage{microtype}
\usepackage{enumitem,xurl,needspace}
\usepackage[dvipsnames]{xcolor}
\usepackage{hyperref}
\hypersetup{unicode=true,colorlinks=true,linkcolor=MidnightBlue,urlcolor=MidnightBlue,
 pdftitle={Research notebook: Problem 12},
 pdfauthor={Open Problems Math research notebook},bookmarksnumbered=true}
\usepackage{bookmark,titlesec,fancyhdr}
\titleformat{\section}{\Large\bfseries\raggedright}{\thesection}{0.7em}{}
\pagestyle{fancy}\fancyhf{}
\fancyhead[L]{\small\sffamily Open Problems Math\enspace|\enspace Research notebook}
\fancyhead[R]{\small\sffamily Problem 12}
\fancyfoot[L]{\small\sffamily 27 September 2026}
\fancyfoot[R]{\small\sffamily \thepage}
\setlength{\parindent}{0pt}
\setlength{\parskip}{5pt plus 1pt}
\setlength{\emergencystretch}{3em}
\setcounter{secnumdepth}{1}
\setcounter{tocdepth}{1}
\setlist[itemize]{leftmargin=3em,itemsep=5pt,topsep=4pt,parsep=0pt}
\allowdisplaybreaks
\newcommand{\N}{\mathbb{N}}
\newcommand{\Z}{\mathbb{Z}}
\newcommand{\Q}{\mathbb{Q}}
\newcommand{\R}{\mathbb{R}}
\newcommand{\C}{\mathbb{C}}
\newcommand{\F}{\mathbb{F}}
\newcommand{\A}{\mathbb{A}}
\newcommand{\PP}{\mathbb P}
\newcommand{\E}{\mathbb{E}}
\newcommand{\eps}{\varepsilon}
\newcommand{\dd}{\,\mathrm d}
\newcommand{\sref}[2]{\hyperlink{source:#1:#2}{[S#2]}}
\newcommand{\eref}[2]{\hyperlink{extra:#1:#2}{[E#2]}}
\newif\ifshowcatalogue
\showcataloguetrue
\newcommand{\cataloguescope}[1]{\ifshowcatalogue
 \paragraph{Exact scope in the supplied catalogue.}
 \begin{quote}\small #1\end{quote}\fi}
\numberwithin{equation}{section}
\begin{document}
\setcounter{section}{0}
% BEGIN PRESERVED SECTION
% ================================================================
% problemId: 12
% Canonical problem ID: 12
% exactTarget SHA-256: 2e57a9156fa8cfdcabacdad86d35b60b3a94d628a9e1288f3b8fe9260471b643
\clearpage
\section*{\texorpdfstring{Goldbach's Conjecture}{Goldbach's Conjecture}}\label{12}
\subsection{Definitions and mathematical statement}
With $\mathcal P$ denoting the positive primes, the statement is
\[
 \forall n\in2\Z,\quad n>2\quad\Longrightarrow\quad
 \exists p,q\in\mathcal P\text{ with }n=p+q.
\]
The primes are not required to be distinct, so $4=2+2$ is allowed. Every even integer is quantified, not merely all sufficiently large integers or a density-one subset.
\par\smallskip\textit{Formulation and concept references:} \sref{12}{1}.
\cataloguescope{Prove or disprove that every even integer n greater than 2 is the sum of two primes.}
\subsection{Short English statement}
Can every even integer larger than two be written as the sum of two primes?
% BEGIN RESEARCH 25
\subsection{Research attempt}

\paragraph{Current frontier.}
The target is every even integer greater than two, not eventual validity
or a density-one result. The published verification through $4\cdot10^{18}$
is a useful finite benchmark, not an argument about unbounded integers.
\eref{12}{1}. Helfgott's theorem for odd integers greater than five is a
three-prime theorem, not the binary assertion. \eref{12}{2}.
Bhowmik--Grimmelt's August 2026 version surveys exceptional-set estimates
and develops an explicit major-arc formula; its additional Siegel-zero
consequence assumes a sparse Hardy--Littlewood estimate. \sref{12}{2}.
Li--Liu's August 2026 revision states that every sufficiently large even
$N$ has $N=p+rq$ with $p,q$ prime, $r=1$ or prime, and $r\le q^{0.9}$.
That recent manuscript theorem, whose complete proof was not independently
audited here, still permits $r>1$ and a finite initial range.
\sref{12}{3}. None of these statements is treated as the exact target.

\paragraph{Attack route.}
I pursue a pointwise convolution lower bound, explicitly accounting for
proper prime powers and the finite initial segment. The missing estimate
is individual control of the minor-arc coefficient, not a mean-square
bound allowing exceptions. A competing idea is that sufficiently small
nonzero Fourier coefficients, together with the correct average density,
might force positive pair counts. An exact quadratic-residue model below
falsifies that proposed inference even with square-root Fourier bounds.
Thus the calculation tests what more arithmetic information the selected
route needs, rather than merely repeating a heuristic independence model.

\paragraph{Attempt.}
Let $\Lambda(m)=\log p$ if $m=p^a$ for a prime $p$ and integer $a\ge1$,
and zero otherwise. All logarithms in this entry are natural, except
$\log_2$. Define
\[
 R(N)=\sum_{m=1}^{N-1}\Lambda(m)\Lambda(N-m),\qquad
 R_{\rm prime}(N)=\sum_{p+q=N}(\log p)(\log q),
\]
where the latter sum is ordered and includes $p=q$. Thus
$R_{\rm prime}(N)>0$ is exactly the desired representation for $N$.

\textbf{Lemma 25.1 (a fully explicit prime-power error).}
For every integer $N\ge4$,
\begin{equation}\label{r25:powers}
 0\le R(N)-R_{\rm prime}(N)
 \le B(N):=2\sqrt N\lfloor\log_2N\rfloor(\log N)^2=o(N).
\end{equation}
\emph{Proof.} For each exponent $a\ge2$ with $2^a\le N$, the number
of powers $p^a\le N$ is at most $N^{1/a}\le\sqrt N$. Summing over at
most $\lfloor\log_2N\rfloor$ exponents bounds the number of proper
prime powers, even with overcounting. In every term of the difference,
one of $m,N-m$ is a proper prime power. There are at most twice that
many ordered terms, and each weight is at most $(\log N)^2$.
Nonnegativity is immediate, and $N^{-1/2}(\log N)^3\to0$ gives the
last statement.\hfill$\square$

Let $e(t)=e^{2\pi it}$ and
$S_N(\alpha)=\sum_{m=1}^{N}\Lambda(m)e(m\alpha)$. Orthogonality of
finite exponential sums gives the exact identity
\begin{equation}\label{r25:circle}
       R(N)=\int_0^1S_N(\alpha)^2e(-N\alpha)\,d\alpha.
\end{equation}
There is no infinite-series interchange here. For a chosen $P\ge1$,
let $\mathfrak M_N$ be the union on $\R/\Z$ of intervals of radius
$P/N$ centered at reduced fractions $a/q$ with $1\le q\le P$;
write $\mathfrak m_N$ for the complement. Define $I_{\mathfrak M}$
and $I_{\mathfrak m}$ by restricting the integral in
\eqref{r25:circle}. Both are real: each domain is invariant under
$\alpha\mapsto-\alpha$, which conjugates the integrand.

The expected major-arc coefficient is
\begin{equation}\label{r25:singular}
 \mathfrak S(N)=2C_2\prod_{\substack{p\mid N\\p>2}}\frac{p-1}{p-2},
 \qquad C_2=\prod_{p>2}\left(1-\frac1{(p-1)^2}\right)>0
 \quad(N\text{ even}).
\end{equation}
The positivity and convergence of $C_2$ follow because every factor is
positive and $\sum_{p>2}(p-1)^{-2}<\infty$; for $0\le u\le1/4$,
$|\log(1-u)|\le(4/3)u$. In particular $\mathfrak S(N)\ge2C_2$.
Its role in the major-arc analysis and the limitations of the error terms
are detailed in \sref{12}{2}, Sections~2--5. Positivity of this formal
main term alone does not prove positivity of the actual coefficient.

\textbf{Lemma 25.2 (a criterion retaining every quantified integer).}
Suppose an explicit $N_0\ge4$ and a constant $c>0$ are supplied such that
for every even $N>N_0$ one has
\begin{equation}\label{r25:missing}
 I_{\mathfrak M}(N)\ge cN,\qquad
 |I_{\mathfrak m}(N)|\le cN/2,\qquad B(N)<cN/2.
\end{equation}
If every even integer $4\le N\le N_0$ is certified as a sum of two
primes, then the full binary Goldbach statement follows.

\emph{Proof.} For $N>N_0$, \eqref{r25:circle} gives $R(N)\ge cN/2$.
Subtracting \eqref{r25:powers} yields $R_{\rm prime}(N)>0$. The finitely
many remaining even integers are covered by the certificates, including
$4=2+2$.\hfill$\square$

The implication is complete, but its analytic hypotheses have not been
proved here. Standard major-arc asymptotics motivate the first condition;
the unresolved universal minor-arc condition is much stronger than an
estimate holding outside a small exceptional set. Nor does the bare
existence of an ineffective threshold eliminate the need to handle its
finite predecessors. An eventual result cannot be presented as a proof
that no small exceptional integer exists.

\medskip
\textbf{Why average control does not force pointwise positivity.}
Suppose an analysis on even $N\in[X,2X]$ gives
$R(N)=\mathfrak S(N)N+E_N$. If $N$ in this range is a counterexample,
then $R(N)\le B(N)$, so for sufficiently large $X$ it satisfies
$|E_N|\ge C_2X$. Consequently
\begin{equation}\label{r25:exceptional}
 \#\{N\in[X,2X]:N\text{ even and exceptional}\}
 \le\frac1{C_2^2X^2}\sum_{\substack{X\le N\le2X\\N\text{ even}}}|E_N|^2.
\end{equation}
An upper bound $O(X^{3-\delta})$ with $0<\delta<1$ on the sum gives
$O(X^{1-\delta})$ exceptions, which still allows one exception or an
infinite sparse sequence. The inference is an exact counting inequality,
not a criticism of the genuine strength of exceptional-set results.
The distinction is central to \sref{12}{2}.

Also, discarding the oscillating factor in the minor-arc integral gives
only
\[
 |I_{\mathfrak m}(N)|\le\int_0^1|S_N(\alpha)|^2d\alpha
                    =\sum_{m\le N}\Lambda(m)^2.
\]
Even the elementary bound $N(\log N)^2$ is too large for
\eqref{r25:missing}; sharper average estimates do not automatically
supply the required individual cancellation. The following finite-group
model makes that logical obstruction particularly explicit.

\textbf{Lemma 25.3 (square-root Fourier cancellation with a missing pair sum).}
Let $p\equiv3\pmod4$ be prime and let $\chi$ be its quadratic character,
with $\chi(0)=0$. On the additive group $\F_p$ define
$g(0)=0$ and $g(x)=1+\chi(x)$ for $x\ne0$. Then
\begin{equation}\label{r25:paley}
 g\ge0,\quad \sum_xg(x)=p-1,\quad (g*g)(0)=0,
 \quad |\widehat g(a)|=\sqrt{p+1}\quad(a\ne0),
\end{equation}
where $\widehat g(a)=\sum_xg(x)e(ax/p)$ and
$(g*g)(z)=\sum_xg(x)g(z-x)$.

\emph{Proof.} Exactly half the nonzero elements are squares, so $g$
is two on them and zero elsewhere. Since $-1$ is not a square,
$\chi(-x)=-\chi(x)$ for $x\ne0$, which gives $g(x)g(-x)=0$,
including at zero. This proves the first three assertions.
For $a\ne0$ put $\tau_a=\sum_x\chi(x)e(ax/p)$. Then
$\widehat g(a)=-1+\tau_a$. Complex conjugation and $x\mapsto-x$
give $\overline{\tau_a}=-\tau_a$. To compute its modulus, substitute
$x=ty$ for $x,y\ne0$:
\[
 |\tau_a|^2
 =\sum_{t\ne0}\chi(t)\sum_{y\ne0}e(ay(t-1)/p)
 =(p-1)-\sum_{t\ne0,1}\chi(t)=p.
\]
The inner sum is $p-1$ for $t=1$ and $-1$ otherwise, while
$\sum_{t\ne0}\chi(t)=0$. Thus $\tau_a$ is purely imaginary of
modulus $\sqrt p$, and $|-1+\tau_a|=\sqrt{p+1}$, proving
\eqref{r25:paley}.\hfill$\square$

The average of $g$ tends to one and every nontrivial Fourier coefficient
is of square-root order, but one convolution coefficient vanishes
exactly. Such primes are arbitrarily large: if only finitely many primes
$3\pmod4$ existed, $4$ times their product minus one would have a new
prime divisor $3\pmod4$. Hence this is not merely a fixed-size anomaly.
In Fourier inversion the nonzero frequencies can collectively cancel
the positive zero-frequency contribution. The calculation is an exact
stress test, not a claim about the distribution of actual primes, and
not a Goldbach counterexample. It shows why a route based only on
average density and generic Fourier cancellation requires further
arithmetic structure in the two-summand case.

\paragraph{Stress test.}
The threshold in Lemma~25.2 is part of its hypothesis. The published
finite verification can close a gap only when an independently justified
analytic threshold lies inside that verified range, not merely because
both statements involve large numbers. Proper prime powers were
subtracted explicitly; positivity of $R(N)$ alone need not yield two
primes. The repeated-prime case is never excluded. No GRH, no
Landau--Siegel-zero exclusion, and no unproved distribution conjecture
has been silently used to establish \eqref{r25:missing}.

Li--Liu's parameter $1.9$ restricts the smaller factor $r$ of the almost
prime $rq$; it does not mean a nonintegral count of prime factors or force
$r=1$. \sref{12}{3}. The quadratic-residue model is not proposed as an
arithmetic substitute for primes: it only falsifies the stated general
Fourier implication. Better structured major/minor-arc arguments may
avoid that obstruction, but must prove their extra estimates.

The companion program uses an exact sieve to produce a prime pair for
every even integer from four through $20{,}000$. It also checks the
proper-prime-power counting bound and the quadratic-character identities
for small primes $3\pmod4$, without numerical Fourier approximations.
These are reproducible finite checks, far below the published range;
they add no evidence of a new analytic threshold and no proof for
unbounded $N$.

\Needspace{8\baselineskip}
\paragraph{Outcome.}
\textbf{Outcome: Conditional reduction.}

The attempt gives an explicit all-integer positivity criterion, with
prime-power contamination and finite-range obligations separated, and
proves a square-root Fourier countermodel to a tempting generic shortcut.
The central uniform minor-arc estimate has not been established, and no
Goldbach counterexample has been found. The classical circle-method
framework and finite-field calculation are not claimed as historically
new. The remaining step is a pointwise arithmetic estimate strong enough
to exclude every exceptional even integer, together with a certified
initial range, not just a smaller exceptional-set density.

% END RESEARCH 25

\subsection{Sources}
\begin{itemize}\raggedright
\item[\textnormal{[S1]}]\hypertarget{source:12:1}{} C. K. Caldwell, "Prime Conjectures and Open Questions," The PrimePages, https://t5k.org/notes/conjectures/.
\url{https://t5k.org/notes/conjectures/}.
{\small\textit{Catalogue formulation source.}}
\item[\textnormal{[S2]}]\hypertarget{source:12:2}{} The exceptional set of the Goldbach problem.
\url{https://arxiv.org/html/2607.27282v2}.
{\small\textit{Catalogue research/status reference; not a proof endorsement.}}
\item[\textnormal{[S3]}]\hypertarget{source:12:3}{} Theorem (1+1.9) on the Goldbach Conjecture.
\url{https://arxiv.org/abs/2606.05224}.
{\small\textit{Catalogue research/status reference; not a proof endorsement.}}
\item[\textnormal{[S4]}]\hypertarget{source:12:4}{} Restricted Goldbach Sums in Arithmetic Progressions: Local Obstructions, an Explicit Almost-All Framework, and a General-Modulus Extension.
\url{https://www.preprints.org/manuscript/202607.2077}.
{\small\textit{Catalogue research/status reference; not a proof endorsement.}}
% BEGIN ADDED REFERENCES 25
\item[\textnormal{[E1]}]\hypertarget{extra:12:1}{} Tom\'as Oliveira e Silva,
Siegfried Herzog and Silvio Pardi, \emph{Empirical verification of the even
Goldbach conjecture and computation of prime gaps up to $4\cdot10^{18}$},
Mathematics of Computation 83 (2014), 2033--2060.
\url{https://doi.org/10.1090/S0025-5718-2013-02787-1}.
Author's project page: \url{https://sweet.ua.pt/tos/goldbach.html}.
{\small\textit{Published finite benchmark and author's account consulted
27 September 2026; not asserted to settle the unbounded problem.}}
\item[\textnormal{[E2]}]\hypertarget{extra:12:2}{} Harald A. Helfgott,
\emph{The ternary Goldbach conjecture is true}, arXiv:1312.7748.
\url{https://arxiv.org/abs/1312.7748}.
{\small\textit{Three-prime theorem and its distinct quantifiers;
abstract consulted 27 September 2026.}}

% END ADDED REFERENCES 25
\end{itemize}

% END PRESERVED SECTION
\end{document}
